\documentclass[12pt]{article}


\usepackage[margin=1in]{geometry}
\usepackage{amsmath,amssymb,amsthm,mathtools}
\usepackage{bm}
\usepackage{mathrsfs}
\usepackage[T1]{fontenc}
\usepackage{lmodern}
\usepackage{microtype}
\usepackage{setspace}
\usepackage{booktabs}
\usepackage{array}
\usepackage{longtable}
\usepackage{caption}
\usepackage{enumitem}
\usepackage{siunitx}
\usepackage{xcolor}
\usepackage[colorlinks=true,allcolors=blue]{hyperref}


\setlength{\parindent}{0pt}
\setlength{\parskip}{0.6em}

\onehalfspacing

\captionsetup{
    font=small,
    labelfont=bf,
    justification=centering
}

\numberwithin{equation}{section}


\DeclareMathOperator{\esssup}{ess\,sup}
\DeclareMathOperator{\Lip}{Lip}
\DeclareMathOperator{\Var}{Var}
\DeclareMathOperator{\diag}{diag}

\newcommand{\R}{\mathbb{R}}
\newcommand{\N}{\mathbb{N}}
\newcommand{\norm}[1]{\left\lVert #1 \right\rVert}
\newcommand{\abs}[1]{\left| #1 \right|}
\newcommand{\Q}{\mathfrak{Q}}
\newcommand{\eps}{\varepsilon}
\newcommand{\dd}{\,\mathrm{d}}


\newtheorem{theorem}{Theorem}[section]
\newtheorem{lemma}[theorem]{Lemma}
\newtheorem{proposition}[theorem]{Proposition}
\newtheorem{corollary}[theorem]{Corollary}

\theoremstyle{definition}

\newtheorem{definition}[theorem]{Definition}
\newtheorem{assumption}[theorem]{Assumption}

\theoremstyle{remark}

\newtheorem{remark}[theorem]{Remark}
\newtheorem{example}[theorem]{Example}


\title{
\textbf{Scott's Coherence Functional}\\[0.5em]
\large A Nonlinear Trajectory-Dependent Framework
for Damping, Oscillatory Memory, and Higher-Order
Dynamical Interactions
}

\author{
Jaxson Scott Levoy Freeman\\
\small Independent student researcher\\
\small Corresponding author:
\texttt{jaxsonstormy@gmail.com}
}

\date{November 2026}


\begin{document}

\maketitle


\begin{abstract}

This paper introduces Scott's Coherence Functional,
a multicomponent mathematical construction designed
to investigate interactions among nonlinear state
dependence, accumulated temporal effects,
exponential attenuation, oscillatory phase,
higher-order derivatives, and acceleration-dependent
contributions.

The proposed functional is formulated through three
principal components: a nonlocal integral contribution
involving mixed second-order spatial derivatives,
a product-weighted third-order derivative contribution,
and a nonlinear correction incorporating velocity
and acceleration.

The construction is examined within an appropriate
space of twice continuously differentiable
trajectories. Under sufficient regularity,
boundedness, and sign assumptions, we establish
well-definedness, continuity in time, and an explicit
global estimate on finite intervals.

We further investigate the mathematical structure
of the integral kernel, derive estimates for
exponentially attenuated memory contributions,
and examine parameter sensitivity with respect
to generalized damping coefficients.

Additional results concern perturbations of
trajectories, limiting regimes, nonlinear
attenuation, and possible extensions to higher-order
dynamical systems.

The analysis also identifies limitations of the
current construction, particularly concerning
trajectory-independent bounds, variational
differentiability, conservation properties, and
numerical convergence.

The purpose of this work is to establish a
mathematically explicit framework for examining
the structural properties of Scott's Coherence
Functional and to provide a foundation for
subsequent analytical investigation.

\end{abstract}

\noindent\textbf{Keywords:}
coherence functional; nonlinear dynamics;
functional analysis; oscillatory integrals;
exponential damping; memory kernels;
trajectory perturbation; asymptotic analysis;
higher-order derivatives.

\medskip

\noindent\textbf{2020 Mathematics Subject Classification:}
Primary 26B15; Secondary 34A34, 45J05, 47H30.

\newpage

\tableofcontents

\newpage


\section{Introduction}

Many mathematical systems involve interactions
between multiple variables that evolve continuously
over time.

Such systems may contain several different
mechanisms simultaneously, including nonlinear
coupling, accumulated phase, exponential
attenuation, state-dependent weighting, and
higher-order derivatives.

A mathematical representation of these
interactions must account for both the instantaneous
configuration of a system and the influence of
its evolution over an interval of time.

The purpose of this paper is to introduce a
trajectory-dependent functional that combines
these mechanisms into a single mathematical
construction.

We refer to this construction as
\textit{Scott's Coherence Functional}.

The functional is designed around three
mathematically distinguishable contributions.

The first is a nonlocal integral term combining
mixed second-order spatial derivatives,
exponential damping, accumulated oscillatory
phase, and nonlinear velocity dependence.

The second introduces third-order derivatives
weighted by products of state-dependent factors.

The third introduces a nonlinear correction
involving both velocity and acceleration.

Together, these contributions provide a
framework for investigating the relationships
between instantaneous state information,
historical trajectory dependence, and
higher-order dynamical quantities.

\subsection{Motivation}

A purely instantaneous mathematical expression
may be insufficient when the behavior of a
system depends on its preceding evolution.

Integral expressions provide one way of
representing such dependence.

In particular, an integral whose upper limit
varies with observation time can incorporate
information from earlier states while allowing
the influence of those states to depend on
subsequent evolution.

Similarly, exponential factors can represent
attenuation, while trigonometric factors can
represent oscillatory interactions.

Nonlinear denominators provide additional
mechanisms for modifying the relative influence
of state variables, velocities, and accelerations.

Scott's Coherence Functional combines these
ideas into a single expression.

\subsection{Principal Mathematical Questions}

The construction raises several questions:

\begin{enumerate}[label=\arabic*.]

\item Under what regularity conditions is the
functional well-defined?

\item Can its magnitude be bounded uniformly
over a finite observation interval?

\item How does accumulated damping influence
historical trajectory contributions?

\item What role does oscillatory phase play
in the resulting integral?

\item Is the functional continuous with respect
to time and perturbations of the trajectory?

\item How sensitive is the construction to
changes in its defining parameters?

\item What limiting behavior occurs under
strong damping, weak damping, and short-time
conditions?

\item Can the construction be generalized to
higher-order or variational mathematical systems?

\end{enumerate}

The following sections develop the mathematical
framework required to investigate these questions.


\section{Definition and Mathematical Framework}

\subsection{State trajectories}

Let \(n\in\N\) be a positive integer and
let \(T>0\) be a finite observation horizon.

Consider a trajectory

\[
\mathbf{x}(t)
=
(x_1(t),x_2(t),\ldots,x_n(t))
\]

taking values in \(\R^n\).

Initially, assume that

\[
\mathbf{x}\in C^2([0,T],\R^n).
\]

The first and second derivatives are denoted by

\[
\dot{x}_i(t)
=
\frac{\dd x_i}{\dd t},
\]

and

\[
\ddot{x}_i(t)
=
\frac{\dd^2x_i}{\dd t^2}.
\]

The functions \(\Phi_i\) and \(\Psi_i\)
describe state-dependent quantities.

The functions \(\omega_i\) and
\(\Omega_{ijk}\) determine damping and
oscillatory contributions.

The function \(\Theta_{ijk}\) represents
an additional accumulated phase contribution.

The coefficients \(\chi_m\), \(\kappa_{ij}\),
and \(\Lambda_{ijk}\) determine the strength
of the corresponding nonlinear terms.

\subsection{Functional spaces}

For a fixed interval \([0,T]\), define

\[
\mathcal{X}_T
=
C^2([0,T],\R^n).
\]

Equip this space with the norm

\[
\begin{aligned}
\norm{\mathbf{x}}_{C^2}
={}&
\max_{t\in[0,T]}
\norm{\mathbf{x}(t)}
\\
&+
\max_{t\in[0,T]}
\norm{\dot{\mathbf{x}}(t)}
\\
&+
\max_{t\in[0,T]}
\norm{\ddot{\mathbf{x}}(t)}.
\end{aligned}
\]

With this norm, \(\mathcal{X}_T\) is a
Banach space.

This provides a natural setting for studying
trajectory-dependent mathematical quantities
because the norm simultaneously controls
the state, velocity, and acceleration.

Assume that

\[
\Phi_i:\R^n\times[0,T]\to\R
\]

and

\[
\Psi_i:\R^n\times[0,T]\to\R.
\]

All spatial derivatives appearing below are
evaluated at the corresponding trajectory
point.

For example,

\[
\frac{\partial^2\Phi_i}
{\partial x_j\partial x_k}
\]

inside the integral denotes

\[
\frac{\partial^2\Phi_i}
{\partial x_j\partial x_k}
(\mathbf{x}(s),s).
\]

Likewise, the mixed derivative in the
second component is evaluated as

\[
\frac{\partial^3\Phi_i}
{\partial x_i^2\partial t}
(\mathbf{x}(t),t),
\]

where the partial derivative with respect
to time is taken while holding the spatial
argument fixed.

\subsection{Definition of the functional}

\begin{definition}[Scott's Coherence Functional]

Let \(T>0\), let \(n\in\N\), and let

\[
\mathbf{x}\in\mathcal{X}_T.
\]

Whenever all required derivatives and
integrals exist, define

\[
\Q:
\mathcal{X}_T\times[0,T]\to\R
\]

by the expression

\[
\begin{aligned}
\Q(\mathbf{x},t)
={}&
\sum_{i,j,k=1}^{n}
\int_0^t
\frac{\partial^2\Phi_i}
{\partial x_j\partial x_k}
(\mathbf{x}(s),s)
\\
&\quad\cdot
\exp\left(
-\int_s^t
\frac{\omega_i(\mathbf{x}(\xi))}
{1+|\dot{x}_j(\xi)|}
\,\dd\xi
\right)
\\
&\quad\cdot
\sin\left(
\int_s^t
\Omega_{ijk}(\mathbf{x}(\eta))
\,\dd\eta
+\Theta_{ijk}(s,t)
\right)
\\
&\quad\cdot
\frac{\dot{x}_j(s)\dot{x}_k(s)}
{1+\displaystyle\sum_{m=1}^{n}
\chi_mx_m^2(s)}
\\
&\quad\cdot
\left(
1+
\frac{
\partial\Psi_i/\partial x_j
(\mathbf{x}(s),s)
}{
1+
\left|
\partial\Psi_k/\partial x_i
(\mathbf{x}(s),s)
\right|
}
\right)
\dd s
\\[8pt]
&+
\sum_{i=1}^{n}
\frac{\partial^3\Phi_i}
{\partial x_i^2\partial t}
(\mathbf{x}(t),t)
\prod_{j=1}^{n}
(1+\kappa_{ij}x_j^2(t))
\\[8pt]
&+
\sum_{i,j,k=1}^{n}
\Lambda_{ijk}
\frac{
\dot{x}_i(t)\dot{x}_j(t)\ddot{x}_k(t)
}{
1+
|\dot{x}_i(t)\dot{x}_j(t)|
+
|\ddot{x}_k(t)|
}.
\end{aligned}
\tag{2.1}
\]

The resulting quantity is called Scott's
Coherence Functional.

\end{definition}

The functional can be decomposed into
three principal components:

\[
\Q
=
\Q_I+\Q_D+\Q_A.
\]

Here:

- \(\Q_I\) denotes the nonlocal integral contribution.
- \(\Q_D\) denotes the product-weighted derivative contribution.
- \(\Q_A\) denotes the nonlinear acceleration contribution.

This decomposition will be used throughout
the subsequent analysis.

\subsection{The integral component}

Define

\[
\begin{aligned}
\Q_I(\mathbf{x},t)
={}&
\sum_{i,j,k=1}^{n}
\int_0^t
H_{ijk}(\mathbf{x},s)
\\
&\qquad\cdot
E_{ij}(s,t)
\sin(P_{ijk}(s,t))
\,\dd s,
\end{aligned}
\]

where the instantaneous trajectory-dependent
factor is

\[
\begin{aligned}
H_{ijk}(\mathbf{x},s)
={}&
\frac{\partial^2\Phi_i}
{\partial x_j\partial x_k}
(\mathbf{x}(s),s)
\\
&\cdot
\frac{\dot{x}_j(s)\dot{x}_k(s)}
{1+\sum_{m=1}^{n}\chi_mx_m^2(s)}
\\
&\cdot
\left(
1+
\frac{
\partial\Psi_i/\partial x_j
(\mathbf{x}(s),s)
}{
1+
\left|
\partial\Psi_k/\partial x_i
(\mathbf{x}(s),s)
\right|
}
\right).
\end{aligned}
\]

The exponential attenuation factor is

\[
E_{ij}(s,t)
=
\exp\left(
-\int_s^t
\frac{\omega_i(\mathbf{x}(\xi))}
{1+|\dot{x}_j(\xi)|}
\,\dd\xi
\right).
\]

The accumulated oscillatory phase is

\[
P_{ijk}(s,t)
=
\int_s^t
\Omega_{ijk}(\mathbf{x}(\eta))
\,\dd\eta
+
\Theta_{ijk}(s,t).
\]

This decomposition separates the
instantaneous trajectory contribution from
the historical weighting mechanism.

\subsection{The derivative component}

The second component is defined by

\[
\Q_D(\mathbf{x},t)
=
\sum_{i=1}^{n}
\frac{\partial^3\Phi_i}
{\partial x_i^2\partial t}
(\mathbf{x}(t),t)
\prod_{j=1}^{n}
(1+\kappa_{ij}x_j^2(t)).
\]

It introduces higher-order derivative
information combined with nonlinear
state-dependent weighting.

\subsection{The acceleration component}

The final component is

\[
\Q_A(\mathbf{x},t)
=
\sum_{i,j,k=1}^{n}
\Lambda_{ijk}
\frac{
\dot{x}_i(t)\dot{x}_j(t)\ddot{x}_k(t)
}{
1+
|\dot{x}_i(t)\dot{x}_j(t)|
+
|\ddot{x}_k(t)|
}.
\]

This component introduces nonlinear
dependence on both velocity and acceleration.


\section{Assumptions and Preliminary Results}

A rigorous analysis requires sufficient
conditions on the coefficient functions,
trajectory, and parameters.

The following assumptions will be used
throughout the principal existence and
boundedness results.

\begin{assumption}[Regularity]

The functions \(\Phi_i\) possess continuous
spatial derivatives through order three.

Their mixed derivatives

\[
\frac{\partial^3\Phi_i}
{\partial x_i^2\partial t}
\]

are continuous.

The functions \(\Psi_i\), \(\omega_i\),
and \(\Omega_{ijk}\) are continuous
wherever evaluated.

The phase functions \(\Theta_{ijk}\)
are continuous on

\[
\Delta_T
=
\{(s,t)\in\R^2:0\leq s\leq t\leq T\}.
\]

\end{assumption}

\begin{assumption}[Nonnegative state weights]

The coefficients satisfy

\[
\chi_m\geq0,
\qquad
\kappa_{ij}\geq0.
\]

All coefficients \(\Lambda_{ijk}\)
are finite real numbers.

\end{assumption}

\begin{assumption}[Nonnegative damping]

For every relevant state,

\[
\omega_i(\mathbf{x})\geq0.
\]

\end{assumption}

\begin{assumption}[Trajectory regularity]

The trajectory satisfies

\[
\mathbf{x}\in C^2([0,T],\R^n).
\]

\end{assumption}

\begin{remark}

These assumptions are sufficient conditions
rather than necessary conditions.

In particular, nonnegative damping guarantees
that the exponential factor does not exceed
one in absolute value.

If negative damping is permitted, the
exponential factor may instead produce
amplification.

\end{remark}

\subsection{Bounded trajectory quantities}

Since the trajectory and its derivatives
are continuous on a compact interval,
the following quantities are finite:

\[
X=
\max_{t\in[0,T]}
\norm{\mathbf{x}(t)},
\]

\[
V=
\max_{\substack{1\leq i\leq n\\t\in[0,T]}}
|\dot{x}_i(t)|,
\]

and

\[
A=
\max_{\substack{1\leq i\leq n\\t\in[0,T]}}
|\ddot{x}_i(t)|.
\]

These constants will be used in the
subsequent estimates.

\subsection{Auxiliary coefficient bounds}

Define

\[
M_2=
\max_{i,j,k}
\sup_{0\leq s\leq T}
\left|
\frac{\partial^2\Phi_i}
{\partial x_j\partial x_k}
(\mathbf{x}(s),s)
\right|.
\]

Similarly, define

\[
P=
\max_{i,j}
\sup_{0\leq s\leq T}
\left|
\frac{\partial\Psi_i}
{\partial x_j}
(\mathbf{x}(s),s)
\right|.
\]

For the third-order derivative contribution,
let

\[
M_3=
\max_i
\sup_{0\leq t\leq T}
\left|
\frac{\partial^3\Phi_i}
{\partial x_i^2\partial t}
(\mathbf{x}(t),t)
\right|.
\]

Finally, define

\[
K=\max_{i,j}\kappa_{ij},
\qquad
L=\max_{i,j,k}|\Lambda_{ijk}|.
\]

All these constants are finite under
the stated assumptions.


\section{Existence and Well-Definedness}

\begin{theorem}[Existence on Finite Time Intervals]

Suppose the regularity, nonnegative-weight,
nonnegative-damping, and trajectory
assumptions hold.

Then Scott's Coherence Functional is finite
and well-defined for every

\[
t\in[0,T].
\]

\end{theorem}

\begin{proof}

Consider first the integral component.

The trajectory is continuous, and the
coefficient functions have the required
continuity properties.

Consequently, the mixed derivative factor

\[
\frac{\partial^2\Phi_i}
{\partial x_j\partial x_k}
(\mathbf{x}(s),s)
\]

is continuous.

The denominator

\[
1+\sum_{m=1}^{n}\chi_mx_m^2(s)
\]

is at least one because

\[
\chi_m\geq0.
\]

Likewise,

\[
1+
\left|
\frac{\partial\Psi_k}{\partial x_i}
\right|
\geq1.
\]

Thus, neither denominator can vanish.

The damping integral is finite because
its integrand is continuous on a compact
interval.

Therefore, the exponential factor is
well-defined and finite.

The phase integral is also finite, and
the sine function is defined for every
real phase argument.

It follows that the complete integrand
is continuous and bounded on the relevant
integration interval.

Its integral therefore exists and is finite.

Next, consider the derivative component.

The mixed third-order derivatives are
continuous, and the trajectory is bounded.

Since the coefficients \(\kappa_{ij}\)
are finite, every factor

\[
1+\kappa_{ij}x_j^2(t)
\]

is finite.

The finite product and finite summation
therefore produce a finite quantity.

Finally, consider the acceleration component.

Its denominator satisfies

\[
1+
|\dot{x}_i\dot{x}_j|
+
|\ddot{x}_k|
\geq1.
\]

The numerator is bounded because velocity
and acceleration are bounded on \([0,T]\).

All coefficients \(\Lambda_{ijk}\) are
finite, and the sums contain finitely
many terms.

Thus, the acceleration component is also
finite.

All three components are therefore
well-defined, proving the theorem.

\end{proof}

\subsection{A bound for the integral component}

The integral component can be bounded
explicitly.

First, nonnegative damping gives

\[
0<E_{ij}(s,t)\leq1.
\]

Furthermore,

\[
|\sin(P_{ijk}(s,t))|\leq1.
\]

The state-dependent denominator satisfies

\[
1+\sum_{m=1}^{n}\chi_mx_m^2(s)\geq1.
\]

The derivative ratio satisfies

\[
\left|
1+
\frac{
\partial\Psi_i/\partial x_j
}{
1+
|\partial\Psi_k/\partial x_i|
}
\right|
\leq1+P.
\]

Consequently,

\[
|H_{ijk}(\mathbf{x},s)|
\leq M_2V^2(1+P).
\]

\begin{proposition}[Integral Component Bound]

Under the preceding assumptions,

\[
|\Q_I(\mathbf{x},t)|
\leq
n^3TM_2V^2(1+P).
\]

\end{proposition}

\begin{proof}

Applying the triangle inequality gives

\[
\begin{aligned}
|\Q_I(\mathbf{x},t)|
\leq
\sum_{i,j,k=1}^{n}
\int_0^t
&
|H_{ijk}(\mathbf{x},s)|
\\
&\cdot E_{ij}(s,t)
|\sin(P_{ijk}(s,t))|
\,\dd s.
\end{aligned}
\]

Using the established estimates,

\[
\begin{aligned}
|\Q_I(\mathbf{x},t)|
&\leq
\sum_{i,j,k=1}^{n}
\int_0^t
M_2V^2(1+P)\,\dd s
\\
&\leq
n^3TM_2V^2(1+P).
\end{aligned}
\]

This proves the proposition.

\end{proof}


\section{Global Bounds for the Functional}

We now establish bounds for the derivative
and acceleration contributions.

\subsection{Derivative component estimate}

Since

\[
|x_j(t)|\leq X
\]

and

\[
\kappa_{ij}\leq K,
\]

we have

\[
1+\kappa_{ij}x_j^2(t)
\leq1+KX^2.
\]

Therefore,

\[
\prod_{j=1}^{n}
(1+\kappa_{ij}x_j^2(t))
\leq
(1+KX^2)^n.
\]

It follows that

\[
|\Q_D(\mathbf{x},t)|
\leq
nM_3(1+KX^2)^n.
\]

\subsection{Acceleration component estimate}

The denominator satisfies

\[
1+
|\dot{x}_i\dot{x}_j|
+
|\ddot{x}_k|
\geq1.
\]

Hence,

\[
\left|
\frac{
\dot{x}_i\dot{x}_j\ddot{x}_k
}{
1+
|\dot{x}_i\dot{x}_j|
+
|\ddot{x}_k|
}
\right|
\leq V^2A.
\]

Since

\[
|\Lambda_{ijk}|\leq L,
\]

we obtain

\[
|\Q_A(\mathbf{x},t)|
\leq n^3LV^2A.
\]

\subsection{Main global boundedness theorem}

\begin{theorem}[Global Bound on a Finite Interval]

Under the assumptions of this paper,

\[
\begin{aligned}
|\Q(\mathbf{x},t)|
\leq{}&
n^3TM_2V^2(1+P)
\\
&+
nM_3(1+KX^2)^n
\\
&+
n^3LV^2A.
\end{aligned}
\]

for every \(t\in[0,T]\).

\end{theorem}

\begin{proof}

The decomposition

\[
\Q=\Q_I+\Q_D+\Q_A
\]

and the triangle inequality yield

\[
|\Q|
\leq
|\Q_I|+|\Q_D|+|\Q_A|.
\]

Substituting the three previously
established component bounds gives

\[
\begin{aligned}
|\Q|
\leq{}&
n^3TM_2V^2(1+P)
\\
&+
nM_3(1+KX^2)^n
\\
&+
n^3LV^2A.
\end{aligned}
\]

This proves the theorem.

\end{proof}

\begin{remark}

The estimate depends on the trajectory
through \(X\), \(V\), and \(A\).

It is therefore a finite-interval estimate
rather than a universal bound independent
of the trajectory.

Additional assumptions would be necessary
to obtain a trajectory-independent estimate.

\end{remark}


\section{The Damping and Oscillatory Components}

\subsection{Exponential damping}

The exponential factor is

\[
E_{ij}(s,t)
=
\exp\left(
-\int_s^t
\frac{\omega_i(\mathbf{x}(\xi))}
{1+|\dot{x}_j(\xi)|}
\,\dd\xi
\right).
\]

This factor introduces accumulated attenuation
over the interval from \(s\) to \(t\).

Because the damping function is nonnegative,

\[
0<E_{ij}(s,t)\leq1.
\]

The velocity-dependent denominator modifies
the effective damping rate.

In particular, define

\[
r_{ij}(\xi)
=
\frac{\omega_i(\mathbf{x}(\xi))}
{1+|\dot{x}_j(\xi)|}.
\]

For fixed nonnegative \(\omega_i\),
increasing the velocity magnitude
decreases this rate.

\begin{proposition}[Damping Monotonicity]

For every \(a\geq0\), the function

\[
f(v)=\frac{a}{1+v},
\qquad v\geq0,
\]

is nonincreasing.

\end{proposition}

\begin{proof}

Differentiation gives

\[
f'(v)
=
-\frac{a}{(1+v)^2}
\leq0.
\]

Therefore, increasing \(v\) cannot increase
the effective damping rate.

\end{proof}

\subsection{Accumulated oscillatory phase}

The oscillatory factor is

\[
\sin\left(
\int_s^t
\Omega_{ijk}(\mathbf{x}(\eta))
\,\dd\eta
+
\Theta_{ijk}(s,t)
\right).
\]

Its argument contains two contributions.

The first is the accumulated trajectory-dependent
phase.

The second is the additional phase
\(\Theta_{ijk}\).

The combination permits oscillatory behavior
to depend on both the evolving state and
an independent phase contribution.

Since

\[
-1\leq\sin(\theta)\leq1,
\]

the oscillatory factor cannot increase
the absolute value of the integrand beyond
the remaining multiplicative factors.

\begin{remark}

The phase contribution can change the sign
of individual portions of the integral.

Consequently, different regions of the
integration interval may reinforce or
partially cancel one another.

This makes the integral component sensitive
to phase relationships along the trajectory.

\end{remark}


\section{Nonlinear State and Acceleration Contributions}

\subsection{State-dependent attenuation}

The integral component contains the factor

\[
R(\mathbf{x},s)
=
\frac{1}
{1+\sum_{m=1}^{n}\chi_mx_m^2(s)}.
\]

Under the nonnegative-weight assumption,

\[
0<R(\mathbf{x},s)\leq1.
\]

For fixed velocities, increasing the weighted
state magnitude decreases this factor.

Thus, the state-dependent denominator
introduces a nonlinear attenuation mechanism.

\subsection{Product-weighted derivative contribution}

The derivative component contains

\[
W_i(\mathbf{x},t)
=
\prod_{j=1}^{n}
(1+\kappa_{ij}x_j^2(t)).
\]

Since every coefficient is nonnegative,

\[
W_i(\mathbf{x},t)\geq1.
\]

Unlike the denominator in the integral
component, this product increases with
the magnitude of the weighted state
coordinates.

It therefore acts as a multiplicative
state-dependent weighting mechanism.

\subsection{Acceleration correction}

Consider

\[
\Q_A
=
\sum_{i,j,k=1}^{n}
\Lambda_{ijk}
\frac{
\dot{x}_i\dot{x}_j\ddot{x}_k
}{
1+
|\dot{x}_i\dot{x}_j|
+
|\ddot{x}_k|
}.
\]

For convenience, define

\[
u=\dot{x}_i\dot{x}_j,
\qquad
v=\ddot{x}_k.
\]

Then

\[
\left|
\frac{uv}{1+|u|+|v|}
\right|
=
\frac{|u||v|}{1+|u|+|v|}.
\]

Because

\[
1+|u|+|v|\geq|u|
\]

and

\[
1+|u|+|v|\geq|v|,
\]

we obtain

\[
\left|
\frac{uv}{1+|u|+|v|}
\right|
\leq
\min(|u|,|v|).
\]

Thus,

\[
\left|
\frac{
\dot{x}_i\dot{x}_j\ddot{x}_k
}{
1+
|\dot{x}_i\dot{x}_j|
+
|\ddot{x}_k|
}
\right|
\leq
\min\left(
|\dot{x}_i\dot{x}_j|,
|\ddot{x}_k|
\right).
\]

This inequality provides an additional
estimate for the acceleration contribution.

\subsection{Structural interpretation}

The acceleration term combines two velocity
factors with one acceleration factor.

Its denominator guarantees that the
expression remains well-defined for
all real velocity and acceleration values.

It also moderates the growth of the
numerator when the quantities become large.

However, the denominator does not imply
that the entire acceleration contribution
is universally bounded independently
of the trajectory.

Such a conclusion would require further
restrictions on the admissible velocities
and accelerations.


\section{Continuity and Time Dependence}

The preceding results establish existence
and boundedness.

We now examine continuity.

\begin{theorem}[Continuity in Time]

Suppose the regularity assumptions hold
and all coefficient functions are continuous
along the trajectory.

Then

\[
t\longmapsto\Q(\mathbf{x},t)
\]

is continuous on \([0,T]\).

\end{theorem}

\begin{proof}

Consider the integral component.

Its integrand is continuous in the variables
\(s\) and \(t\) on the compact triangular
domain

\[
\Delta_T
=
\{(s,t):0\leq s\leq t\leq T\}.
\]

The exponential and sine functions are
continuous compositions of the accumulated
integrals.

All denominators remain strictly positive.

The parameter-dependent integral therefore
varies continuously with its upper limit
and with its integrand.

Hence, every summand in the integral
component is continuous in \(t\).

The derivative component is continuous
because its derivatives and state-dependent
factors are continuous.

The acceleration component is continuous
because the trajectory belongs to \(C^2\)
and its denominator is strictly positive.

Since the functional is a finite sum of
these continuous components, it is
continuous on \([0,T]\).

\end{proof}

\begin{corollary}[Attainment of Extreme Values]

The functional attains a maximum and
a minimum on every compact interval
\([0,T]\).

\end{corollary}

\begin{proof}

The functional is continuous on the
compact interval \([0,T]\).

The extreme value theorem therefore
guarantees that it attains both a maximum
and a minimum.

\end{proof}


\section{Nonlocal Structure and Memory Dependence}

An important feature of the integral
component is that its value at time \(t\)
depends on trajectory information accumulated
over the interval preceding \(t\).

This property distinguishes the integral
component from purely instantaneous
nonlinear expressions.

\subsection{Construction of the memory kernel}

Recall the definitions

\[
D_{ij}(s,t)
=
\int_s^t
\frac{\omega_i(\mathbf{x}(\xi))}
{1+|\dot{x}_j(\xi)|}
\,\dd\xi
\]

and

\[
P_{ijk}(s,t)
=
\int_s^t
\Omega_{ijk}(\mathbf{x}(\eta))
\,\dd\eta
+
\Theta_{ijk}(s,t).
\]

Define the combined memory kernel by

\[
K_{ijk}(s,t)
=
e^{-D_{ij}(s,t)}
\sin(P_{ijk}(s,t)).
\]

The integral component can consequently
be represented as

\[
\Q_I(\mathbf{x},t)
=
\sum_{i,j,k=1}^{n}
\int_0^t
K_{ijk}(s,t)
H_{ijk}(\mathbf{x},s)
\,\dd s.
\]

This representation separates the historical
weighting mechanism from the instantaneous
trajectory-dependent contribution.

\subsection{Kernel magnitude estimate}

Under nonnegative damping,

\[
D_{ij}(s,t)\geq0.
\]

Therefore,

\[
\begin{aligned}
|K_{ijk}(s,t)|
&=
e^{-D_{ij}(s,t)}
|\sin(P_{ijk}(s,t))|
\\
&\leq e^{-D_{ij}(s,t)}
\\
&\leq1.
\end{aligned}
\]

Thus,

\[
|K_{ijk}(s,t)|\leq1.
\]

\subsection{Exponential memory attenuation}

Suppose there exists a constant
\(\delta>0\) such that

\[
\frac{\omega_i(\mathbf{x}(\xi))}
{1+|\dot{x}_j(\xi)|}
\geq\delta
\]

throughout the relevant interval.

Then

\[
D_{ij}(s,t)
\geq\delta(t-s).
\]

Consequently,

\[
|K_{ijk}(s,t)|
\leq e^{-\delta(t-s)}.
\]

This establishes an exponentially decaying
upper bound on the influence of earlier
trajectory states.

\subsection{Uniform memory-weight estimate}

Integrating the preceding inequality gives

\[
\int_0^t
|K_{ijk}(s,t)|\,\dd s
\leq
\int_0^t
e^{-\delta(t-s)}\,\dd s.
\]

Evaluating the integral,

\[
\begin{aligned}
\int_0^t
e^{-\delta(t-s)}\,\dd s
&=
\frac{1-e^{-\delta t}}{\delta}.
\end{aligned}
\]

Therefore,

\[
\int_0^t
|K_{ijk}(s,t)|\,\dd s
\leq\frac{1}{\delta}.
\]

The bound is independent of the observation
time \(t\).

This demonstrates that a uniformly positive
effective damping rate produces a uniformly
bounded total memory weight, even as the
observation interval increases.

\begin{remark}

The estimate controls the total magnitude
of the memory kernel.

It does not imply that every individual
historical contribution is positive, since
the oscillatory phase may change its sign.

\end{remark}


\section{Parameter Sensitivity and Perturbation Analysis}

The behavior of Scott's Coherence Functional
depends on several independent parameters.

These include damping coefficients,
oscillatory coefficients, nonlinear state
weights, and acceleration coefficients.

A systematic analysis of their influence
provides additional information concerning
the structural sensitivity of the construction.

\subsection{Generalized damping parameter}

Introduce a nonnegative parameter
\(\alpha\) into the damping expression:

\[
E_{ij}^{(\alpha)}(s,t)
=
\exp(-\alpha D_{ij}(s,t)).
\]

For \(\alpha=1\), the original damping
factor is recovered.

Differentiating with respect to \(\alpha\),

\[
\frac{\partial E_{ij}^{(\alpha)}}
{\partial\alpha}
=
-D_{ij}(s,t)
e^{-\alpha D_{ij}(s,t)}.
\]

Since \(D_{ij}(s,t)\geq0\),

\[
\frac{\partial E_{ij}^{(\alpha)}}
{\partial\alpha}
\leq0.
\]

Thus, increasing the damping parameter
cannot increase the exponential attenuation
factor.

A second differentiation yields

\[
\frac{\partial^2 E_{ij}^{(\alpha)}}
{\partial\alpha^2}
=
D_{ij}^2(s,t)
e^{-\alpha D_{ij}(s,t)}
\geq0.
\]

The attenuation factor is therefore
decreasing and convex as a function
of \(\alpha\).

\subsection{Sensitivity of the integral component}

Let \(\Q_I^{(\alpha)}\) denote the integral
component with generalized damping.

Under the established regularity conditions,
differentiation under the integral sign gives

\[
\begin{aligned}
\frac{\partial\Q_I^{(\alpha)}}{\partial\alpha}
=
-\sum_{i,j,k=1}^{n}
\int_0^t
&D_{ij}(s,t)
e^{-\alpha D_{ij}(s,t)}
\\
&\cdot
\sin(P_{ijk}(s,t))
H_{ijk}(\mathbf{x},s)
\,\dd s.
\end{aligned}
\]

This derivative measures the sensitivity
of the integral component to changes
in the damping parameter.

Importantly, its sign is not necessarily
negative because the oscillatory factor
and the trajectory-dependent contribution
may change sign.

Therefore, although the exponential
attenuation decreases monotonically
with \(\alpha\), the complete integral
component need not do so.

\subsection{Sensitivity bounds}

Since

\[
D e^{-\alpha D}\leq\frac{1}{e\alpha}
\]

for \(D\geq0\) and \(\alpha>0\), we obtain

\[
\begin{aligned}
\left|
\frac{\partial\Q_I^{(\alpha)}}{\partial\alpha}
\right|
&\leq
\frac{n^3tM_2V^2(1+P)}
{e\alpha}.
\end{aligned}
\]

Hence,

\[
\left|
\frac{\partial\Q_I^{(\alpha)}}{\partial\alpha}
\right|
\leq
\frac{n^3TM_2V^2(1+P)}
{e\alpha}.
\]

This provides a quantitative estimate
of parameter sensitivity for positive
damping parameters.

\subsection{Perturbations of the trajectory}

Consider two trajectories

\[
\mathbf{x},\mathbf{y}
\in\mathcal{X}_T.
\]

Define their separation by

\[
\eps
=
\norm{\mathbf{x}-\mathbf{y}}_{C^2}.
\]

A natural question is whether small
changes in the trajectory produce small
changes in the functional.

Assume additionally that the coefficient
functions and all derivatives required
by the functional are locally Lipschitz
on a common bounded region containing
the images of both trajectories.

\begin{theorem}[Local Perturbation Stability]

Under the preceding additional assumptions,
there exists a finite constant \(C_T\)
such that

\[
\sup_{t\in[0,T]}
|\Q(\mathbf{x},t)-\Q(\mathbf{y},t)|
\leq
C_T
\norm{\mathbf{x}-\mathbf{y}}_{C^2}.
\]

\end{theorem}

\begin{proof}

The functional is composed of finitely many
products, quotients, compositions, and
integrals of trajectory-dependent functions.

On a common bounded region, the relevant
coefficient functions and their derivatives
are locally Lipschitz.

Furthermore, all denominators are bounded
away from zero by the nonnegative-weight
assumptions.

Consequently, the instantaneous factors
in each summand satisfy local Lipschitz
estimates with respect to the state,
velocity, and acceleration.

The accumulated damping and phase terms
inherit corresponding estimates because
their integration intervals have uniformly
bounded length.

The exponential and sine functions are
Lipschitz on bounded intervals.

Applying the product and composition
estimates to the integral, derivative,
and acceleration components yields
finite constants \(C_I,C_D,C_A\) such that

\[
\sup_t
|\Q_I(\mathbf{x},t)-\Q_I(\mathbf{y},t)|
\leq C_I\eps,
\]

\[
\sup_t
|\Q_D(\mathbf{x},t)-\Q_D(\mathbf{y},t)|
\leq C_D\eps,
\]

and

\[
\sup_t
|\Q_A(\mathbf{x},t)-\Q_A(\mathbf{y},t)|
\leq C_A\eps.
\]

Combining the estimates gives

\[
\sup_t
|\Q(\mathbf{x},t)-\Q(\mathbf{y},t)|
\leq
(C_I+C_D+C_A)\eps.
\]

Setting

\[
C_T=C_I+C_D+C_A
\]

proves the result.

\end{proof}

\subsection{Interpretation of perturbation stability}

The preceding estimate establishes a
distinction between boundedness and stability.

Boundedness controls the magnitude of
the functional along an individual trajectory.

Perturbation stability controls the difference
between functional values associated with
two nearby trajectories.

This distinction is relevant when studying
approximate trajectories, numerical solutions,
and parameter-dependent models.


\section{Asymptotic Behavior and Limiting Cases}

The behavior of Scott's Coherence Functional
under limiting regimes provides additional
information about its mathematical structure.

We consider several cases in which individual
contributions become dominant, negligible,
or asymptotically simplified.

\subsection{The small-time regime}

Consider the limit

\[
t\to0^+.
\]

From the integral bound,

\[
|\Q_I(\mathbf{x},t)|
\leq
n^3tM_2V^2(1+P).
\]

Therefore,

\[
\lim_{t\to0^+}
\Q_I(\mathbf{x},t)=0.
\]

Since the remaining two components are
continuous in time,

\[
\lim_{t\to0^+}
\Q(\mathbf{x},t)
=
\Q_D(\mathbf{x},0)
+
\Q_A(\mathbf{x},0).
\]

Thus, the integral contribution becomes
negligible over sufficiently short
observation intervals.

\subsection{The strong-damping regime}

Suppose the effective damping rate satisfies

\[
\frac{\omega_i(\mathbf{x}(\xi))}
{1+|\dot{x}_j(\xi)|}
\geq\delta
\]

for some \(\delta>0\).

Then

\[
|E_{ij}(s,t)|
\leq e^{-\delta(t-s)}.
\]

For an individual summand, suppose

\[
|H_{ijk}(\mathbf{x},s)|
\leq M_H
\]

for a finite constant \(M_H\).

It follows that

\[
\begin{aligned}
\left|
\int_0^t
E_{ij}(s,t)
\sin\bigl(P_{ijk}(s,t)\bigr)
H_{ijk}(\mathbf{x},s)
\,ds
\right|
&\leq
M_H\int_0^t e^{-\delta(t-s)}\,ds
\\[6pt]
&=
\frac{M_H}{\delta}
\left(1-e^{-\delta t}\right)
\\[6pt]
&\leq
\frac{M_H}{\delta}.
\end{aligned}
\]

Hence, the individual memory contributions
remain uniformly controlled under sufficiently
strong damping.

\subsection{The weak-damping regime}

Suppose the damping coefficient is scaled
by a parameter \(\alpha\) satisfying

\[
\alpha\to0^+.
\]

Then

\[
E_{ij}^{(\alpha)}(s,t)
=
e^{-\alpha D_{ij}(s,t)}
\longrightarrow1.
\]

For a fixed finite observation interval,
the convergence is uniform because the
accumulated damping quantity is bounded.

Consequently, the integral component
approaches the undamped expression

\[
\begin{aligned}
\Q_I^{(0)}(\mathbf{x},t)
=
\sum_{i,j,k=1}^{n}
\int_0^t
&H_{ijk}(\mathbf{x},s)
\\
&\cdot
\sin(P_{ijk}(s,t))
\,\dd s.
\end{aligned}
\]

This limiting case isolates the influence
of accumulated phase and nonlinear trajectory
dependence from exponential attenuation.

\subsection{Large-state behavior}

Consider the state-dependent factor

\[
R(\mathbf{x},s)
=
\frac{1}
{1+\sum_{m=1}^{n}\chi_mx_m^2(s)}.
\]

If at least one positively weighted
state coordinate becomes unbounded along
a sequence, then

\[
R(\mathbf{x},s)\to0.
\]

For bounded velocities and bounded
remaining integrand factors, the corresponding
velocity contribution therefore approaches zero.

This identifies a conditional attenuation
mechanism associated with large weighted
state magnitudes.

\subsection{Limiting behavior of the acceleration term}

Consider the scalar expression

\[
F(u,v)
=
\frac{uv}{1+|u|+|v|}.
\]

If \(u\) is fixed and \(v\to0\), then

\[
\lim_{v\to0}F(u,v)=0.
\]

Likewise, if \(v\) is fixed and \(u\to0\),

\[
\lim_{u\to0}F(u,v)=0.
\]

If both quantities become large with
comparable magnitudes, the expression
can also become large.

Thus, the nonlinear denominator moderates
the growth of the acceleration contribution
but does not establish universal boundedness.


\section{Special Cases and Simplifications}

The general expression contains several
mechanisms that can be examined independently
through particular parameter choices.

\subsection{Vanishing acceleration contribution}

If

\[
\Lambda_{ijk}=0
\]

for every triple of indices, then

\[
\Q_A=0.
\]

The functional reduces to

\[
\Q=\Q_I+\Q_D.
\]

\subsection{Vanishing state weights}

If

\[
\chi_m=0
\]

for every \(m\), then

\[
1+\sum_{m=1}^{n}\chi_mx_m^2=1.
\]

The velocity contribution simplifies to

\[
\dot{x}_j(s)\dot{x}_k(s).
\]

\subsection{Vanishing oscillatory phase}

If

\[
\Omega_{ijk}=0
\]

and

\[
\Theta_{ijk}=0,
\]

then

\[
\sin(0)=0.
\]

Consequently,

\[
\Q_I=0.
\]

The functional then consists of its
derivative and acceleration contributions.

\subsection{Constant trajectories}

Suppose

\[
\mathbf{x}(t)=\mathbf{x}_0
\]

for every \(t\).

Then

\[
\dot{x}_i(t)=0
\]

and

\[
\ddot{x}_i(t)=0.
\]

The integral contribution vanishes because
its numerator contains a product of velocities.

The acceleration contribution also vanishes.

Therefore,

\[
\begin{aligned}
\Q(\mathbf{x}_0,t)
=
\sum_{i=1}^{n}
&
\frac{\partial^3\Phi_i}
{\partial x_i^2\partial t}
(\mathbf{x}_0,t)
\\
&\cdot
\prod_{j=1}^{n}
(1+\kappa_{ij}x_{0,j}^2).
\end{aligned}
\]

This special case isolates the derivative
contribution of the functional.

\subsection{One-dimensional reduction}

For \(n=1\), all triple summations
collapse into a single term.

Writing \(x=x_1\), the functional becomes

\[
\begin{aligned}
\Q(x,t)
={}&
\int_0^t
\Phi_{1,xx}(x(s),s)
E(s,t)
\\
&\cdot
\sin(P(s,t))
\frac{\dot{x}^2(s)}
{1+\chi_1x^2(s)}
\\
&\cdot
\left(
1+
\frac{\Psi_{1,x}(x(s),s)}
{1+|\Psi_{1,x}(x(s),s)|}
\right)
\,\dd s
\\
&+
\Phi_{1,xxt}(x(t),t)
(1+\kappa_{11}x^2(t))
\\
&+
\Lambda_{111}
\frac{\dot{x}^2(t)\ddot{x}(t)}
{1+\dot{x}^2(t)+|\ddot{x}(t)|}.
\end{aligned}
\]

This reduction provides a simpler setting
for numerical experimentation and
analytical examination.


\section{Possible Generalizations}

The functional can be generalized in several
directions without changing its fundamental
three-component structure.

\subsection{Generalized damping}

Introduce a positive parameter \(\alpha\):

\[
E_{ij}^{(\alpha)}(s,t)
=
\exp\left(
-\alpha
\int_s^t
\frac{\omega_i(\mathbf{x}(\xi))}
{1+|\dot{x}_j(\xi)|}
\,\dd\xi
\right).
\]

The parameter controls the strength
of the damping contribution.

For \(\alpha=1\), the original expression
is recovered.

\subsection{Generalized oscillatory functions}

The sine function may be replaced by
a more general bounded continuous function
\(g:\R\to\R\).

The corresponding factor becomes

\[
g(P_{ijk}(s,t)).
\]

If

\[
\sup_{\theta\in\R}|g(\theta)|<\infty,
\]

then the principal existence and
boundedness arguments remain applicable.

\subsection{Higher-order trajectory dependence}

The acceleration term may be extended
to include third-order derivatives.

For example,

\[
\Q_J
=
\sum_{i,j,k=1}^{n}
\Gamma_{ijk}
\frac{
\dot{x}_i\dot{x}_j\dddot{x}_k
}{
1+
|\dot{x}_i\dot{x}_j|
+
|\dddot{x}_k|
}.
\]

Such an extension requires trajectories
with sufficient additional regularity,
for example

\[
\mathbf{x}\in C^3([0,T],\R^n).
\]

\subsection{Alternative nonlinear attenuation}

The state-dependent denominator may
be generalized to

\[
1+
\sum_{m=1}^{n}
\chi_m|x_m|^p,
\qquad p>0.
\]

For nonnegative coefficients, this
expression remains positive.

The parameter \(p\) controls the rate
at which the attenuation changes with
state magnitude.

\subsection{Generalized product weighting}

The derivative contribution may also
be extended through a positive weighting
function \(w_{ij}\):

\[
\Q_D^{(w)}
=
\sum_{i=1}^{n}
\frac{\partial^3\Phi_i}
{\partial x_i^2\partial t}
\prod_{j=1}^{n}
w_{ij}(x_j).
\]

Suitable assumptions on \(w_{ij}\)
would preserve the existence of this
component.

These generalizations suggest that the
original construction can be regarded
as one member of a broader family of
trajectory-dependent nonlinear functionals.


\section{Limitations and Open Mathematical Problems}

Although the preceding analysis establishes
several fundamental properties of Scott's
Coherence Functional, a number of mathematical
questions remain unresolved.

The distinction between established results
and potential extensions is important
for understanding the actual scope
of the construction.

\subsection{Trajectory-independent bounds}

The global estimate established in this
paper depends on

\[
X,\qquad V,\qquad A.
\]

These quantities represent bounds on
the trajectory, velocity, and acceleration.

Consequently, the existing estimate does
not establish a universal upper bound
over the entire space \(\mathcal{X}_T\).

A natural extension would investigate
whether additional restrictions on admissible
trajectories could produce estimates
independent of individual trajectory
magnitudes.

\subsection{Differentiability with respect
to trajectories}

Continuity with respect to trajectories
does not automatically imply differentiability.

A further investigation would examine
whether the mapping

\[
\mathbf{x}\mapsto\Q(\mathbf{x},t)
\]

possesses a Fréchet derivative on appropriate
subsets of the trajectory space.

If such a derivative exists, it would
provide a linear approximation of the
functional under small perturbations.

Specifically, one would seek a bounded
linear operator \(D\Q(\mathbf{x})\)
satisfying

\[
\begin{aligned}
\Q(\mathbf{x}+\mathbf{h},t)
={}&
\Q(\mathbf{x},t)
\\
&+
D\Q(\mathbf{x})[\mathbf{h}]
+
o(\norm{\mathbf{h}}_{C^2}),
\end{aligned}
\]

as

\[
\norm{\mathbf{h}}_{C^2}\to0.
\]

Establishing this property would provide
a foundation for a more complete
variational analysis.

\subsection{Existence of conservation laws}

The current formulation does not establish
that the functional is conserved along
solutions of a particular differential equation.

A conservation property would require
additional structural assumptions.

One possible objective would be to identify
conditions under which

\[
\frac{\dd}{\dd t}
\Q(\mathbf{x},t)=0
\]

along a specified class of trajectories.

Such a result would require an appropriate
evolution equation and a detailed analysis
of the time derivative of each component.

\subsection{Monotonicity and sign structure}

The oscillatory integral component is not
necessarily nonnegative.

Indeed, the sine factor can change sign,
while the derivative and acceleration
contributions may also be positive or negative.

Therefore, the functional cannot generally
be interpreted as a positive-definite quantity.

An interesting open problem is to identify
additional restrictions on the coefficient
functions and trajectories that guarantee
monotonicity or sign definiteness.

\subsection{Numerical approximation}

Another direction concerns the numerical
evaluation of the functional.

The nested integrals and nonlinear
derivative terms introduce several
potential sources of numerical error.

A numerical investigation could compare
quadrature methods and trajectory
approximation schemes.

One possible objective would be to establish
an estimate of the form

\[
|\Q(\mathbf{x},t)-\Q_h(\mathbf{x},t)|
\leq Ch^p,
\]

where \(h\) denotes a discretization
parameter and \(p\) represents the
convergence order.

The derivation of such an estimate would
require assumptions on the smoothness
of the integrand and the selected
numerical approximation method.

\subsection{Existence of stationary trajectories}

Another possible direction is the
identification of trajectories for which
the functional becomes stationary.

A corresponding investigation could seek
trajectories satisfying

\[
D\Q(\mathbf{x})[\mathbf{h}]=0
\]

for every admissible perturbation
\(\mathbf{h}\).

This would require an appropriate
differentiability theory and additional
conditions on the trajectory space.

These unresolved questions provide
several directions for subsequent
mathematical investigation.


\section{Conclusion}

This paper introduced Scott's Coherence
Functional, a trajectory-dependent
mathematical construction combining
integral, derivative, oscillatory,
damping, and acceleration contributions.

The functional contains three principal
components.

The first combines mixed second-order
spatial derivatives with accumulated
exponential damping, oscillatory phase,
and nonlinear velocity dependence.

The second introduces third-order
derivatives weighted by products
of state-dependent factors.

The third introduces a nonlinear
correction involving velocity
and acceleration.

Under appropriate regularity, boundedness,
and sign assumptions, the functional
was shown to be well-defined on finite
time intervals.

An explicit global bound was established
in terms of the trajectory, velocity,
acceleration, and relevant coefficient
functions.

Continuity in time was also demonstrated,
allowing the functional to attain maximum
and minimum values on compact intervals.

The analysis further identified a
nonlocal memory structure within the
integral component.

Under a uniformly positive effective
damping rate, the total absolute memory
weight admits a uniform bound independent
of the observation time.

Parameter sensitivity was examined
through a generalized damping parameter,
and additional assumptions were used
to establish local perturbation stability
with respect to trajectory changes.

Several limiting cases were also
investigated, including short-time,
strong-damping, weak-damping, and
large-state regimes.

The construction admits potential
generalizations involving alternative
oscillatory functions, higher-order
trajectory derivatives, and nonlinear
attenuation mechanisms.

Nevertheless, several mathematical
questions remain open.

These include Fréchet differentiability,
trajectory-independent bounds,
conservation properties, sign definiteness,
and numerical convergence.

A particularly interesting direction
is the development of a complete
variational theory for \(\Q\), examining
how changes in the trajectory influence
the resulting functional.

The present work provides an initial
analytical foundation for these
investigations and establishes a
framework within which the structural
properties of Scott's Coherence Functional
can be studied more systematically.


\section{Acknowledgments}

The author thanks the teachers and mentors
who encouraged independent mathematical
exploration and the development of
original mathematical ideas.

The author also acknowledges the
foundational contributions of researchers
in functional analysis, differential
equations, and mathematical analysis,
whose established results provide
important tools for the present work.


\begin{thebibliography}{99}

\bibitem{conway}
J. B. Conway,

\emph{A Course in Functional Analysis},

2nd ed., Graduate Texts in Mathematics,
vol. 96, Springer, New York, 1990.

\bibitem{rudin}
W. Rudin,

\emph{Principles of Mathematical Analysis},

3rd ed., International Series in Pure
and Applied Mathematics,
McGraw-Hill, New York, 1976.

\bibitem{kreyszig}
E. Kreyszig,

\emph{Introductory Functional Analysis
with Applications},

John Wiley \& Sons, New York, 1978.

\bibitem{walter}
W. Walter,

\emph{Ordinary Differential Equations},

Springer-Verlag, New York, 1998.

\bibitem{evans}
L. C. Evans,

\emph{Partial Differential Equations},

2nd ed., Graduate Studies in Mathematics,
vol. 19, American Mathematical Society,
Providence, 2010.

\bibitem{dieudonne}
J. Dieudonn\'e,

\emph{Foundations of Modern Analysis},

Academic Press, New York, 1960.

\bibitem{cartan}
H. Cartan,

\emph{Differential Calculus},

Hermann, Paris, 1971.

\bibitem{kreyszig2}
E. Kreyszig,

\emph{Advanced Engineering Mathematics},

10th ed., John Wiley \& Sons, 2011.

\end{thebibliography}


\end{document}
